% META: {"id": "1NSI-APT-CO-103", "chapitre": "1NSI-ALGO-PARCOURS-TRIS", "type_objet": "corrige", "exercice_ref": "1NSI-APT-EX-103", "capacites": ["P-ALGO-02A"], "capacites_codes": ["C3"], "sources_inspiration": [], "status": "needs_review"}
% BEGIN-VERIFY
% def tri_insertion(tableau):
%     tableau = tableau[:]
%     for i in range(1, len(tableau)):
%         valeur = tableau[i]
%         j = i - 1
%         while j >= 0 and tableau[j] > valeur:
%             tableau[j + 1] = tableau[j]
%             j = j - 1
%         tableau[j + 1] = valeur
%     return tableau
% assert tri_insertion([5, 2, 9, 1, 7]) == [1, 2, 5, 7, 9]
% assert tri_insertion([]) == []
% assert tri_insertion([3]) == [3]
% assert tri_insertion([2, 2, 1]) == [1, 2, 2]
% def decalages(tableau):
%     tableau = tableau[:]
%     total = 0
%     for i in range(1, len(tableau)):
%         valeur = tableau[i]
%         j = i - 1
%         while j >= 0 and tableau[j] > valeur:
%             tableau[j + 1] = tableau[j]
%             j = j - 1
%             total += 1
%         tableau[j + 1] = valeur
%     return total
% assert decalages([5, 2, 9, 1, 7]) == 5
% assert decalages([1, 2, 5, 7, 9]) == 0
% assert decalages([9, 7, 5, 2, 1]) == 10
% assert 10 == 4 * 5 // 2
% END-VERIFY
\begin{corrige}{1NSI-APT-EX-103}

\textbf{1.} Tour par tour :

\begin{itemize}
  \item $i=1$ : valeur $2$ ; $1$ décalage ($5$ passe à droite) ; tableau
    \lstinline{[2, 5, 9, 1, 7]}.
  \item $i=2$ : valeur $9$ ; $0$ décalage ; tableau \lstinline{[2, 5, 9, 1, 7]}.
  \item $i=3$ : valeur $1$ ; $3$ décalages ($9$, $5$ et $2$ passent à droite) ; tableau
    \lstinline{[1, 2, 5, 9, 7]}.
  \item $i=4$ : valeur $7$ ; $1$ décalage ($9$ passe à droite) ; tableau
    \lstinline{[1, 2, 5, 7, 9]}.
\end{itemize}

\textbf{2.} Le total est de $1+0+3+1 = 5$ décalages. Le tour $i=2$ n'en effectue aucun :
la valeur mise de côté est $9$, et l'élément immédiatement à sa gauche vaut $5$. La
condition \lstinline{tableau[j] > valeur} est fausse dès le premier test, la boucle
\lstinline{while} ne s'exécute pas, et $9$ est réécrit à sa propre place. Un élément déjà
plus grand que toute la zone triée reste où il est.

\textbf{3.}
\begin{python}
def tri_insertion(tableau):
    tableau = tableau[:]
    for i in range(1, len(tableau)):
        valeur = tableau[i]
        j = i - 1
        while j >= 0 and tableau[j] > valeur:
            tableau[j + 1] = tableau[j]
            j = j - 1
        tableau[j + 1] = valeur
    return tableau
\end{python}
\lstinline{tri_insertion([5, 2, 9, 1, 7])} renvoie \lstinline{[1, 2, 5, 7, 9]}. Le test
\lstinline{j >= 0} doit précéder \lstinline{tableau[j] > valeur} : Python évalue de gauche
à droite et s'arrête dès que la première condition est fausse, ce qui évite de lire
\lstinline{tableau[-1]}.

\textbf{4.} Sur \lstinline{[1, 2, 5, 7, 9]} déjà trié : $0$ décalage. Chaque valeur est
supérieure ou égale à celle qui la précède, donc la boucle \lstinline{while} s'arrête
immédiatement à chaque tour.

Sur \lstinline{[9, 7, 5, 2, 1]} trié à l'envers : $10$ décalages. Chaque nouvelle valeur
est plus petite que toutes les précédentes et doit traverser toute la zone triée, soit
$1+2+3+4 = 10$ décalages.

Le coût du tri par insertion dépend donc \textbf{des données}, et pas seulement de leur
nombre : il est très rapide sur un tableau presque trié, et le plus lent sur un tableau
trié à l'envers. C'est ce qui le distingue du tri par sélection, dont le nombre de
comparaisons est le même quelles que soient les valeurs.

\end{corrige}
